\chapter{Inability to Urinate (Urinary Retention)}

\section*{Definition}
Urinary retention is the inability to empty the bladder completely or at all. Acute urinary retention is sudden and painful, while chronic retention develops gradually with incomplete emptying. It can affect both men and women but is most common in older men.

\section*{What Happens in Your Body}
Urinary retention can be obstructive (physical blockage of urine flow) or non-obstructive (failure of detrusor contraction or coordination). Bladder outlet obstruction: prostate enlargement, urethral stricture, pelvic organ prolapse. Detrusor underactivity: neurologic conditions, medications, aging, or damage to bladder innervation. In acute retention, the bladder overdistends, causing severe pain and triggering afferent nerve signals.

\section*{Causes}
\begin{itemize}
\item Benign prostatic hyperplasia (BPH) is the most common cause in men
\item Urethral stricture
\item Prostate cancer
\item Pelvic organ prolapse in women
\item Bladder neck contracture
\item Neurogenic bladder: spinal cord injury, multiple sclerosis, Parkinson disease, diabetic neuropathy
\item Medications: anticholinergics, antihistamines, decongestants, opioids, alpha-agonists
\item Post-surgical: anesthesia, pelvic surgery, post-hip replacement
\item Constipation (fecal impaction compressing urethra)
\item Infection: prostatitis, cystitis, urethritis
\item Anatomical: phimosis, paraphimosis
\end{itemize}

\section*{When to See a Doctor}
\begin{itemize}
\item Inability to urinate or difficulty starting urination
\item Painful, urgent need to urinate but unable to pass urine
\item Lower abdominal pain and distension
\item Feeling of bladder fullness
\item Overflow incontinence (dribbling)
\item Frequency or urgency with small volumes
\item Palpable distended bladder on examination
\end{itemize}

\section*{Related Symptoms}
\begin{itemize}
\item Enlarged prostate (BPH)
\item Frequent urination
\item Overactive bladder
\item Urinary tract infection
\item Hematuria (blood in urine)
\item Constipation
\end{itemize}

\section*{Self-Care / Home Management}
\begin{itemize}
\item Do NOT delay seeking medical care for acute retention
\item Try to relax and find a comfortable position to urinate
\item Run water or place a hand in warm water to stimulate urination
\item Gentle suprapubic pressure (Cred maneuver) may help
\item Avoid alcohol, caffeine, and decongestants
\item Maintain regular bowel movements to avoid constipation
\end{itemize}

\section*{How Your Doctor Will Evaluate You}
\begin{itemize}
\item History and physical examination including digital rectal exam
\item Bladder scan or ultrasound for post-void residual
\item Urinalysis and urine culture
\item Uroflowmetry and pressure-flow studies
\item Cystoscopy to evaluate for obstruction
\item Urodynamic studies for suspected neurogenic bladder
\end{itemize}

\section*{Treatment Options}
\begin{itemize}
\item Immediate catheterization for acute retention (Foley or intermittent)
\item Alpha-blockers (tamsulosin, terazosin) for BPH
\item 5-alpha reductase inhibitors for BPH
\item Clean intermittent catheterization for chronic retention
\item Surgery: TURP, urethrotomy, bladder neck incision
\item Treatment of underlying cause (medication adjustment, infection treatment)
\item Neuromodulation for non-obstructive retention
\end{itemize}

\section*{References}
\begin{thebibliography}{99}
\bibitem{ref1} Selius BA, Subedi R. "Urinary retention in adults: diagnosis and initial management." \textit{American Family Physician}, 2008;77(5):643--650.
\bibitem{ref2} McVary KT, Roehrborn CG, Avins AL, et al. "Update on AUA guideline on the management of benign prostatic hyperplasia." \textit{Journal of Urology}, 2011;185(5):1793--1803.
\bibitem{ref3} Abrams P, Cardozo L, Fall M, et al. "The standardisation of terminology in lower urinary tract function." \textit{European Urology}, 2003;43(1):1--10.
\bibitem{ref4} Kaplan SA, Wein AJ, Staskin DR, et al. "Urinary retention: current concepts." \textit{Urology}, 2019;132:21--28.
\bibitem{ref5} Oelke M, Bachmann A, Descazeaud A, et al. "EAU guidelines on the treatment and follow-up of non-neurogenic male lower urinary tract symptoms." \textit{European Urology}, 2013;64(1):118--140.
\bibitem{ref6} Yoshimura N, Chancellor MB. "Physiology and pharmacology of the bladder." UpToDate, 2023.

\end{thebibliography}

\clearpage
